% =====================================================================
%  Thesis summary (max 2 pages) - Segreteria Carriere Scienze, UniVR
%  Compile: pdflatex Thesis_Summary_Fikadu.tex   (or paste into Overleaf)
% =====================================================================
\documentclass[11pt,a4paper]{article}

\usepackage[T1]{fontenc}
\usepackage[utf8]{inputenc}
\usepackage{lmodern}
\usepackage[english]{babel}
\usepackage{microtype}
\usepackage[top=20mm,bottom=20mm,left=22mm,right=22mm]{geometry}
\usepackage{tabularx}
\usepackage{titlesec}
\usepackage{parskip}

\pdfmapfile{}                        % avoid font-expansion errors on some installs
\microtypesetup{expansion=false}

\titleformat{\section}{\normalfont\bfseries\scshape\normalsize}{}{0pt}{}
\titlespacing*{\section}{0pt}{10pt}{3pt}

\pagestyle{empty}
\setlength{\emergencystretch}{2em}

\newcommand{\metarow}[2]{\textbf{#1} & #2 \\}

\begin{document}

% ------------------------------------------------------------------ header
\begin{center}
  {\large\bfseries\scshape Thesis Summary}
\end{center}
\vspace{4pt}

\noindent
\begin{tabularx}{\textwidth}{@{}p{0.30\textwidth}X@{}}
  \metarow{Surname, Name}{FIKADU, Dagmawi Tarekegn}
  \metarow{Student number}{VR509938}
  \metarow{Degree course}{Master's Degree in Molecular and Medical Biotechnology}
  \metarow{Department}{Department of Biotechnology, University of Verona}
  \metarow{Thesis title}{Curation and Cheminformatic Characterization of a
        Clinical-Stage GPCR Ligand Dataset for Machine-Learning-Based Binding
        Prediction}
  \metarow{Principal Investigator}{Prof.\ Alejandro Giorgetti}
  \metarow{Supervisor}{Filippo Baldessari}
  \metarow{Academic year}{2025--2026}
\end{tabularx}

\vspace{6pt}

% ------------------------------------------------------------------ body
\section{Background and aim}
G protein-coupled receptors (GPCRs) are the largest family of membrane receptors in the
human genome and the target of roughly a third of approved drugs. Predicting how tightly a
small molecule binds a given receptor, from chemical structure alone, would allow large
compound libraries to be prioritised computationally before any experiment is performed.
The obstacle is not the modelling itself but the evaluation: chemical datasets contain
dense series of close analogues, so a model tested on randomly selected compounds is
largely re-predicting molecules it has effectively already seen, and the resulting
performance is inflated. This thesis had two aims: to assemble and characterise a curated
dataset of clinical-stage GPCR ligands, and to build binding-affinity models on that
dataset under validation strict enough that the reported numbers can be trusted.

\section{Dataset curation and characterisation}
Human $K_i$ binding data were retrieved from ChEMBL for compounds in clinical phases
I--IV, restricted to molecules below 700~Da, and processed through a curation chain
comprising salt removal, peptide exclusion, a flexibility filter and deduplication
retaining the highest $\mathrm{pK_i}$ per molecule. The resulting set contains 3,011
compound--target records spanning 719 unique molecules, 47 targets and 18 receptor
families, with a strong aminergic composition. Fifteen molecular descriptors were computed
for every compound.

Two properties of the set proved consequential. First, selecting on clinical status rather
than single-target potency yields extensively polypharmacological chemistry: molecules act
on 4.19 distinct targets on average, and 64.7\% engage more than one. Second, ligand
efficiency declines sharply with molecular size (Spearman $\rho=-0.77$): median efficiency
falls from 0.52 to 0.24 between the smallest and largest compounds while median potency is
essentially unchanged (6.66 against 6.51), so additional molecular mass purchases no
additional affinity. Principal component analysis showed the descriptor space to be
low-dimensional, with 65.5\% of variance in two components, and clinical phase did not
separate meaningfully within it.

\section{Predictive modelling and validation}
Two data-rich aminergic receptors, serotonin 5-HT$_{2A}$ and dopamine D$_3$, were selected
for modelling. The curated clinical set provides only 172 and 173 compounds per target, too
few to train and validate honestly, so each was extended with preclinical ChEMBL compounds
under identical curation rules, giving 4,620 and 5,488 molecules. Each molecule was
represented as a 2,248-dimensional vector combining a 2,048-bit ECFP4 fingerprint with
approximately 200 normalised physicochemical descriptors, and three tree ensembles were
trained.

Validation used a repeated Bemis--Murcko scaffold split in which whole scaffold families
are assigned to one side of the partition, so that every test compound is structurally
novel to the model. Under this protocol the best model reached a mean held-out $R^2$ of
0.562 for 5-HT$_{2A}$ and 0.644 for D$_3$ over twelve partitions, with root-mean-squared
errors of 0.81 and 0.76 log units. Under conventional ten-fold cross-validation the same
models score 0.679 and 0.717, and both estimates are reported so that the gap between them
is explicit. The scaffold-split error corresponds to approximately 1.6 and 1.5 times the
inter-laboratory reproducibility of pooled ChEMBL $K_i$ measurements, indicating that a
substantial share of the residual error originates in the labels rather than the model.

Recast as a virtual-screening classifier at a 100~nM potency cut-off, the models achieved
Matthews correlation coefficients of 0.635 and 0.632, with balanced errors and calibrated
output probabilities, and top-1\% enrichment at approximately 98\% of the ceiling imposed
by the class balance. Four independent controls support these figures: a 100-permutation
y-randomisation test ($p=0.0099$), a nearest-neighbour read-across baseline that the models
exceed by roughly 0.2~$R^2$, an ablation removing polypharmacological compounds, and a
stricter similarity-cluster split. Under that last and most demanding partition performance
approximately halves, and the model retains an advantage over the baseline for D$_3$ but
not for 5-HT$_{2A}$ --- a limit that is reported explicitly. SHAP attribution indicates
that the two receptors are governed by different property families, charge and
electrostatics for 5-HT$_{2A}$ and molecular topology for D$_3$.

\section{Conclusions}
The work delivers a curated and characterised clinical-stage GPCR ligand dataset, and
affinity models whose performance is established under leakage-controlled evaluation and
supported by four separate controls. The principal contribution is methodological: the
reported figures are conservative estimates of prospective performance rather than the
optimistic values that random-split evaluation produces, and the conditions under which the
models fail are stated alongside those under which they succeed. Extension across the wider
receptor panel, three-dimensional and structure-aware representations, and the inclusion of
confirmed inactive compounds are the natural continuations.

% ------------------------------------------------------------------ signature
\vspace{10pt}
\hrule
\vspace{8pt}

\noindent Verona, 14 September 2026

\vspace{22pt}
\noindent
\begin{tabularx}{\textwidth}{@{}p{0.45\textwidth}X@{}}
  \hrulefill & \\[2pt]
  {\footnotesize Signature of the candidate} &
  {\footnotesize Dagmawi Tarekegn Fikadu --- VR509938} \\
\end{tabularx}

\end{document}
